%%%%%%%%%%%%%%%%%%%%%%%%%%%%%%%%%%%%%%%%%
% Research Poster - 3 Column Layout
% Deep Food Image Classifier with Evolutionary HPO
%
% Author: V. Harsha Vardhan Yellela
% Course: MCS-5993 Evolutionary Computing & Deep Learning
% Institution: Lawrence Technological University
%%%%%%%%%%%%%%%%%%%%%%%%%%%%%%%%%%%%%%%%%

\documentclass[a0,landscape]{a0poster}

\usepackage{multicol}
\usepackage[svgnames]{xcolor}
\usepackage{times}
\usepackage{graphicx}
\graphicspath{{figures/}}
\usepackage{booktabs}
\usepackage[font=small,labelfont=bf]{caption}
\usepackage{amsfonts, amsmath, amsthm, amssymb}
\usepackage{wrapfig}
\usepackage{tcolorbox}
\usepackage{enumitem}
\usepackage{geometry}

% Page setup
\geometry{left=1.5cm, right=1.5cm, top=1cm, bottom=1cm}

% Colors
\definecolor{headerblue}{RGB}{0, 84, 147}
\definecolor{boxblue}{RGB}{0, 84, 147}
\definecolor{lightblue}{RGB}{220, 235, 250}

% Custom box style
\tcbset{
    boxstyle/.style={
        colback=white,
        colframe=boxblue,
        fonttitle=\bfseries\large\color{white},
        coltitle=white,
        colbacktitle=boxblue,
        boxrule=2pt,
        arc=3pt,
        left=5pt,
        right=5pt,
        top=5pt,
        bottom=5pt,
        toptitle=3pt,
        bottomtitle=3pt,
    }
}

\begin{document}

%----------------------------------------------------------------------------------------
%	HEADER
%----------------------------------------------------------------------------------------
\begin{minipage}[c]{0.15\linewidth}
\includegraphics[width=\linewidth]{LTU_logo.png}
\end{minipage}
\hfill
\begin{minipage}[c]{0.65\linewidth}
\centering
{\Huge\bfseries\color{headerblue} Deep Food Image Classifier: Comparing\\[0.3cm]
Evolutionary Hyperparameter Optimization Methods}\\[0.5cm]
{\Large V Harsha Vardhan Yellela, MCS Candidate; CJ Chung, PhD, Professor}\\[0.2cm]
{\large College of Arts \& Science, Lawrence Technological University}
\end{minipage}
\hfill
\begin{minipage}[c]{0.15\linewidth}
\includegraphics[width=\linewidth]{LTU_seal.png}
\end{minipage}

\vspace{0.8cm}

%----------------------------------------------------------------------------------------
%	THREE COLUMN LAYOUT
%----------------------------------------------------------------------------------------
\begin{minipage}[t]{0.32\linewidth}

%----------------------------------------------------------------------------------------
% LEFT COLUMN
%----------------------------------------------------------------------------------------

\begin{tcolorbox}[boxstyle, title=ABSTRACT]
This work presents a comprehensive comparison of three hyperparameter optimization (HPO) methods for deep learning: hand-coded ES(1+1) Evolution Strategy, Keras Tuner Hyperband, and DEAP library-based ($\mu$+$\lambda$)-ES. We develop a 10-class American food image classifier using a custom ResNet-style CNN architecture with residual connections, achieving up to \textbf{79.8\% test accuracy}—a \textbf{24\% improvement} over the baseline. Our experiments demonstrate that evolutionary strategies can match or exceed gradient-free hyperparameter search methods while providing interpretable optimization trajectories. The system utilizes multi-GPU mixed-precision training on dual RTX 3090 GPUs, processes 10,000 images across 10 food categories, and optimizes 5 hyperparameters: learning rate, dropout rate, label smoothing, rotation range, and zoom range.
\end{tcolorbox}

\vspace{0.5cm}

\begin{tcolorbox}[boxstyle, title=INTRODUCTION]
Hyperparameter optimization remains a critical challenge in deep learning, significantly impacting model performance. While grid search and random search are common approaches, evolutionary algorithms offer a principled way to explore the hyperparameter space efficiently.

\vspace{0.3cm}
\begin{center}
\includegraphics[width=0.95\linewidth]{food_samples.png}
\captionof{figure}{Sample images from 10 American food classes: chicken wings, churros, french fries, hamburger, hot dog, ice cream, mac \& cheese, pancakes, pizza, waffles}
\end{center}

\textbf{Research Questions:}
\begin{enumerate}[leftmargin=*,noitemsep]
    \item Can hand-coded ES(1+1) match library-based HPO methods?
    \item How do different evolutionary strategies compare for CNN optimization?
    \item What hyperparameters have the most impact on food classification?
\end{enumerate}
\end{tcolorbox}

\vspace{0.5cm}

\begin{tcolorbox}[boxstyle, title=DATASET]
\textbf{Source:} Food-101 dataset (subset of 10 American food classes)

\begin{minipage}[t]{0.48\linewidth}
\textbf{Statistics:}
\begin{itemize}[leftmargin=*,noitemsep]
    \item Total: 10,000 images
    \item 1,000 per class
    \item Train/Val/Test: 70/15/15\%
    \item Image Size: 224×224
\end{itemize}
\end{minipage}
\hfill
\begin{minipage}[t]{0.48\linewidth}
\textbf{Classes:}
\begin{itemize}[leftmargin=*,noitemsep]
    \item Chicken Wings
    \item Churros
    \item French Fries
    \item Hamburger
    \item Hot Dog
\end{itemize}
\end{minipage}

\begin{minipage}[t]{0.48\linewidth}
\textbf{Data Augmentation:}
\begin{itemize}[leftmargin=*,noitemsep]
    \item Random horizontal flip
    \item Random rotation
    \item Random zoom
    \item Random contrast
\end{itemize}
\end{minipage}
\hfill
\begin{minipage}[t]{0.48\linewidth}
\textbf{Classes (cont.):}
\begin{itemize}[leftmargin=*,noitemsep]
    \item Ice Cream
    \item Mac \& Cheese
    \item Pancakes
    \item Pizza
    \item Waffles
\end{itemize}
\end{minipage}
\end{tcolorbox}

\vspace{0.5cm}

\begin{tcolorbox}[boxstyle, title=HPO METHODS COMPARED]
\textbf{1. Baseline CNN:}
\begin{itemize}[leftmargin=*,noitemsep]
    \item Fixed hyperparameters (LR=0.002, Dropout=0.25)
    \item 50 epochs with early stopping
    \item Reference for comparison
\end{itemize}

\textbf{2. ES(1+1) - Hand-coded:}
\begin{itemize}[leftmargin=*,noitemsep]
    \item Single parent, single offspring strategy
    \item Gaussian mutation with 1/5th success rule
    \item $\sigma$ adaptation: $\sigma \leftarrow \sigma \cdot e^{\pm 0.2}$
    \item 15 generations, 30 epochs per evaluation
\end{itemize}

\textbf{3. Keras Tuner (Hyperband):}
\begin{itemize}[leftmargin=*,noitemsep]
    \item Successive halving algorithm
    \item Factor=3, Hyperband iterations=2
    \item 174 trials completed
\end{itemize}

\textbf{4. DEAP ($\mu$+$\lambda$)-ES:}
\begin{itemize}[leftmargin=*,noitemsep]
    \item Population: $\mu$=5, $\lambda$=10
    \item Tournament selection (size=3)
    \item 10 generations
\end{itemize}
\end{tcolorbox}

\end{minipage}
\hfill
%----------------------------------------------------------------------------------------
% MIDDLE COLUMN
%----------------------------------------------------------------------------------------
\begin{minipage}[t]{0.32\linewidth}

\begin{tcolorbox}[boxstyle, title=CNN ARCHITECTURE]
\textbf{Model:} Custom ResNet-style with 4 residual blocks

\begin{center}
\includegraphics[width=0.95\linewidth]{cnn_architecture.png}
\captionof{figure}{ResNet-style CNN Architecture with Residual Connections}
\end{center}

\begin{itemize}[leftmargin=*,noitemsep]
    \item Initial Conv: 64 filters, 7×7 kernel, stride 2
    \item Residual Block 1: 64 → 128 channels
    \item Residual Block 2: 128 → 256 channels
    \item Residual Block 3: 256 → 512 channels
    \item Residual Block 4: 512 → 1024 channels
    \item Global Average Pooling + Dense(1024) + Softmax(10)
    \item \textbf{Total Parameters: $\sim$20.6M}
\end{itemize}
\end{tcolorbox}

\vspace{0.5cm}

\begin{tcolorbox}[boxstyle, title=HYPERPARAMETER OPTIMIZATION]
\textbf{Parameters Optimized:}
\begin{center}
\begin{tabular}{|l|c|c|}
\hline
\textbf{Parameter} & \textbf{Range} & \textbf{Type} \\
\hline
Learning Rate & $[10^{-5}, 10^{-2}]$ & Log \\
Dropout Rate & $[0.1, 0.5]$ & Linear \\
Label Smoothing & $[0.0, 0.3]$ & Linear \\
Rotation Range & $[0.0, 0.2]$ & Linear \\
Zoom Range & $[0.0, 0.3]$ & Linear \\
\hline
\end{tabular}
\end{center}

\vspace{0.3cm}
\begin{center}
\includegraphics[width=0.95\linewidth]{es_convergence.png}
\captionof{figure}{ES(1+1) Fitness Convergence over 15 Generations}
\end{center}
\end{tcolorbox}

\vspace{0.5cm}

\begin{tcolorbox}[boxstyle, title=RESULTS]
\textbf{Test Accuracy Comparison:}
\begin{center}
\begin{tabular}{|l|c|c|}
\hline
\textbf{Method} & \textbf{Test Acc.} & \textbf{Improv.} \\
\hline
Baseline CNN & 64.37\% & --- \\
\hline
\textbf{ES(1+1)} & \textbf{79.80\%} & \textbf{+24.0\%} \\
\hline
Keras Tuner & 79.27\% & +23.1\% \\
\hline
DEAP ($\mu$+$\lambda$) & 79.33\% & +23.2\% \\
\hline
\end{tabular}
\end{center}

\vspace{0.3cm}
\begin{center}
\includegraphics[width=0.95\linewidth]{accuracy_comparison.png}
\captionof{figure}{Test Accuracy Comparison: All Four Methods}
\end{center}
\end{tcolorbox}

\vspace{0.5cm}

\begin{tcolorbox}[boxstyle, title=OPTIMIZED HYPERPARAMETERS]
\begin{center}
\begin{tabular}{|l|c|c|c|c|}
\hline
\textbf{HP} & \textbf{Base} & \textbf{ES} & \textbf{Tuner} & \textbf{DEAP} \\
\hline
LR & 2e-3 & 1.4e-3 & 1.4e-3 & 1e-4 \\
\hline
Dropout & 0.25 & 0.26 & 0.30 & 0.10 \\
\hline
Label Sm. & 0.05 & 0.30 & 0.28 & 0.08 \\
\hline
Rotation & 0.05 & 0.07 & 0.05 & 0.03 \\
\hline
Zoom & 0.10 & 0.11 & 0.10 & 0.33 \\
\hline
\end{tabular}
\end{center}

\begin{center}
\includegraphics[width=0.95\linewidth]{hyperparameter_comparison.png}
\captionof{figure}{Optimized Hyperparameters by Method}
\end{center}
\end{tcolorbox}

\end{minipage}
\hfill
%----------------------------------------------------------------------------------------
% RIGHT COLUMN
%----------------------------------------------------------------------------------------
\begin{minipage}[t]{0.32\linewidth}

\begin{tcolorbox}[boxstyle, title=KEY FINDINGS]
\begin{itemize}[leftmargin=*]
    \item \textbf{All three HPO methods achieved similar performance} ($\sim$79\%)
    \item \textbf{Hand-coded ES(1+1) slightly outperformed} library methods
    \item \textbf{Label smoothing had significant impact} (0.05 → 0.28-0.30)
    \item \textbf{Learning rate convergence:} $\sim$1.4×10$^{-3}$ optimal
    \item \textbf{24\% improvement} over baseline validates HPO importance
\end{itemize}
\end{tcolorbox}

\vspace{0.5cm}

\begin{tcolorbox}[boxstyle, title=TRAINING CONFIGURATION]
\textbf{Hardware:}
\begin{itemize}[leftmargin=*,noitemsep]
    \item 2× NVIDIA RTX 3090 (48GB VRAM total)
    \item 21GB System RAM + 8GB Swap
    \item WSL2 Ubuntu environment
\end{itemize}

\textbf{Software:}
\begin{itemize}[leftmargin=*,noitemsep]
    \item TensorFlow 2.10 with Keras
    \item Mixed Precision (float16) training
    \item MirroredStrategy for multi-GPU
\end{itemize}

\textbf{Training Time:}
\begin{itemize}[leftmargin=*,noitemsep]
    \item Baseline: $\sim$1.5 hours (50 epochs)
    \item ES(1+1): $\sim$6-8 hours (15 gen × 30 epochs)
    \item Keras Tuner: $\sim$6 hours (174 trials)
    \item DEAP: $\sim$4 hours (10 gen × 30 epochs)
\end{itemize}

\begin{center}
\includegraphics[width=0.95\linewidth]{training_curves.png}
\captionof{figure}{Training and Validation Accuracy Curves}
\end{center}
\end{tcolorbox}

\vspace{0.5cm}

\begin{tcolorbox}[boxstyle, title=INFERENCE \& DEPLOYMENT]
\textbf{Out-of-Scope Detection:}
\begin{itemize}[leftmargin=*,noitemsep]
    \item Confidence threshold: 60\%
    \item Predictions below threshold flagged as "out of scope"
    \item Prevents misclassification of non-food images
\end{itemize}

\textbf{Model Deployment:}
\begin{itemize}[leftmargin=*,noitemsep]
    \item Hugging Face Hub: \texttt{HAR5HA-YELLELA/american\_food\_classifier}
    \item Inference time: 3-5 seconds per image
    \item Memory: 4-5 GB VRAM for inference
\end{itemize}

\begin{center}
\includegraphics[width=0.95\linewidth]{inference_example.png}
\captionof{figure}{Inference Example: Churros classified with 75.1\% confidence}
\end{center}
\end{tcolorbox}

\vspace{0.5cm}

\begin{tcolorbox}[boxstyle, title=CONCLUSION]
\begin{itemize}[leftmargin=*,noitemsep]
    \item All HPO methods effective: 23-24\% improvement over baseline
    \item ES(1+1) competitive: Hand-coded implementation matched library methods
    \item Label smoothing critical: Increased from 0.05 to 0.28-0.30
    \item Evolutionary strategies viable: For CNN hyperparameter optimization
    \item Practical deployment: Model available on Hugging Face Hub
\end{itemize}

\textbf{Recommendations:}
\begin{enumerate}[leftmargin=*,noitemsep]
    \item Start with ES(1+1) for simplicity and interpretability
    \item Use Keras Tuner for extensive search with limited manual coding
    \item Consider DEAP for complex multi-objective optimization
\end{enumerate}
\end{tcolorbox}

\vspace{0.5cm}

\begin{tcolorbox}[boxstyle, title=FUTURE WORK]
\begin{enumerate}[leftmargin=*,noitemsep]
    \item Extend to full Food-101 dataset (101 classes)
    \item Implement CMA-ES for covariance matrix adaptation
    \item Multi-objective optimization (accuracy vs. inference time)
    \item Neural Architecture Search (NAS) with evolutionary methods
    \item Mobile deployment with model quantization
\end{enumerate}
\end{tcolorbox}

\vspace{0.5cm}

\begin{tcolorbox}[boxstyle, title=REFERENCES]
\small
\begin{enumerate}[leftmargin=*,noitemsep]
    \item Bossek, J. et al. (2019). "A Survey on Hyperparameter Optimization." \textit{arXiv:1907.10902}
    \item Li, L. et al. (2018). "Hyperband: A Novel Bandit-Based Approach." \textit{JMLR}
    \item Fortin, F.A. et al. (2012). "DEAP: Evolutionary Algorithms Made Easy." \textit{JMLR}
    \item He, K. et al. (2016). "Deep Residual Learning." \textit{CVPR}
\end{enumerate}
\end{tcolorbox}

\end{minipage}

\end{document}
